\subsection{The identity}

BBQ defines its ambiguous-context bias score as
$s_{\text{AMB}} = (1 - \text{acc}) \cdot s_{\text{DIS}}$
\citep{parrish2022bbq} (top row of Figure~\ref{fig:method}), motivating the scaling as reflecting that a biased
answer ``is more harmful if it happens more often.'' On ambiguous items the
gold answer \emph{is} the \textsc{unknown} option, so $\text{acc}$ is
precisely the abstention rate $p$, and writing $c = 1-p$ for the rate at
which the model commits to a group answer, and $b$ for the directional bias
among committed answers,
\[
s_{\text{AMB}} = c \cdot b .
\]
Any reported change therefore decomposes \emph{exactly}, with no modeling
assumptions and no access to per-item data:
\[
\underbrace{s_1 - s_0}_{\text{reported}} =
\underbrace{b_0\,(c_1 - c_0)}_{\text{coverage}} +
\underbrace{c_1\,(b_1 - b_0)}_{\text{disposition}} .
\]
The second term captures a change in committed-answer bias, which may
reflect both changes in responses and changes in the composition of answered
items. The allocation between terms follows a decomposition convention and is not a
unique causal attribution (symmetric alternative in Appendix~\ref{app:decomp}). The first term says the intervention moved
mass into the abstain option, widening the Manski identified set for the
underlying preference from width $p_0$ to $p_1$. Holding committed-answer bias fixed, increasing abstention moves
the score toward zero while widening the identified interval. This
establishes a mechanical link, not an equivalence between every score
improvement and a loss of identification.
